\chapter{Literature Survey}
This chapter presents a comprehensive review of existing research in Retrieval-Augmented Generation (RAG), syntactic ambiguity in Natural Language Processing (NLP), and the application of quantum computing in language understanding. The objective is to analyze the limitations of current systems and establish the foundation for the proposed hybrid QRAG framework.

\section{Retrieval-Augmented Generation}

Retrieval-Augmented Generation (RAG) was introduced to overcome the limitations of parametric knowledge in Large Language Models (LLMs) by integrating them with external knowledge sources \cite{lewis2021rag}. This paradigm enhances factual accuracy by retrieving relevant documents and incorporating them into the generation process.

However, the performance of RAG systems is highly dependent on the quality of retrieved context. Studies have shown that models struggle to effectively utilize long-context information, especially when dealing with complex queries \cite{liu2023lostmiddle}. Additionally, while vector-based methods show promise in detecting semantic similarity, they tend to fail at identifying explicit relations, which has led to an increasing trend toward dynamic knowledge graphs. It is evident from existing literature that there is a major limitation in automatically generating graph structures; in particular, models often hallucinate or skip non-conformant entities, leading to sparse data structures. Enhancements such as iterative calibration and concept-based retrieval have been proposed to improve retrieval alignment and contextual relevance \cite{sgic2025,concept2025rag}.

\section{Syntactic Ambiguity in NLP}

Syntactic ambiguity remains a major challenge in NLP, where multiple interpretations arise due to structural variations in sentences. Classical parsing techniques often rely on local heuristics, leading to incorrect interpretations.

Research indicates that models exhibit local attachment bias, associating modifiers with nearby tokens instead of considering global structure \cite{trustworthy2025}. Additionally, modern LLMs demonstrate overconfidence in ambiguous scenarios, often producing incorrect yet fluent outputs \cite{racquet2025}. These limitations highlight the need for more expressive models capable of capturing structural dependencies.

\section{Quantum Machine Learning}

Quantum Machine Learning (QML) introduces a novel approach to handling complex data by leveraging high-dimensional Hilbert spaces. Quantum-enhanced feature representations allow for improved modeling of non-linear relationships \cite{havlicek2019quantum}.

Variational Quantum Classifiers (VQC) provide a practical implementation of QML using parameterized quantum circuits, enabling learning in complex feature spaces \cite{schuld2020circuit,schuld2021encoding}. These models are particularly suitable for near-term quantum devices operating in the NISQ era.

\section{Quantum Natural Language Processing}

Quantum Natural Language Processing (QNLP) extends quantum principles to linguistic modeling by representing words as quantum states and grammar as quantum operations.

The DisCoCat framework establishes a connection between grammatical structures and quantum circuits, enabling compositional representation of sentence meaning \cite{coeke2020qnlp}. Tools such as lambeq facilitate the transformation of classical parse trees into executable quantum circuits \cite{kartsaklis2021lambeq}. These approaches enable the modeling of non-local dependencies in language.

\section{Hybrid Quantum-Classical Systems}

Due to the limitations of current quantum hardware, hybrid quantum-classical systems have emerged as a practical solution. These systems leverage classical computation for efficiency while utilizing quantum models for complex tasks \cite{preskill2018nisq}.

Research suggests that quantum advantage is task-dependent and becomes significant only for problems involving high complexity and non-linear relationships \cite{schuld2021encoding}. This observation motivates the need for selective integration of quantum components in NLP systems.

\section{Comparative Analysis of Existing Works}

To better understand the strengths and limitations of existing research, a comparative analysis of key studies is presented in Table~\ref{tab:literature_comparison}.

\begin{table}[h]
    \centering
    \caption{Comparison of Existing Research Works}
    \label{tab:literature_comparison}
    \begin{tabular}{|p{3cm}|p{3cm}|p{4cm}|p{4cm}|}
        \hline
        \textbf{Work} & \textbf{Domain} & \textbf{Key Contribution} & \textbf{Limitation} \\
        \hline
        Lewis et al. (2021) \cite{lewis2021rag} & RAG & Introduced retrieval-augmented generation for improving factual accuracy & Fails in handling syntactic ambiguity \\
        \hline
        Liu et al. (2023) \cite{liu2023lostmiddle} & RAG & Identified long-context utilization issues in LLMs & Poor handling of complex contextual dependencies \\
        \hline
        Havlíček et al. (2019) \cite{havlicek2019quantum} & QML & Proposed quantum-enhanced feature spaces for classification & Requires specialized quantum hardware \\
        \hline
        Schuld et al. (2020) \cite{schuld2020circuit} & QML & Introduced circuit-based quantum classifiers & Limited scalability on NISQ devices \\
        \hline
        Coecke et al. (2020) \cite{coeke2020qnlp} & QNLP & Developed DisCoCat framework for quantum linguistic modeling & High implementation complexity \\
        \hline
        Kartsaklis et al. (2021) \cite{kartsaklis2021lambeq} & QNLP & Provided toolkit for mapping NLP to quantum circuits & Limited real-world deployment \\
        \hline
        Preskill (2018) \cite{preskill2018nisq} & Quantum Computing & Defined NISQ-era limitations and hybrid approaches & Noise and hardware constraints persist \\
        \hline
    \end{tabular}
\end{table}

\section{Research Gap}

As is evident from the literature, there exists a notable weakness in the current approaches. While large RAGs have improved factual consistency and can be further optimized using algorithmic verification, they are unable to effectively reduce the deep structural ambiguity caused by inherent local attachment biases. On the other hand, quantum approaches offer greater representational strength but are currently not being applied to any NLP system.

To address this issue, a new approach involving the selective use of quantum computation for linguistics in classical systems needs to be explored.

\bigskip\noindent
From the literature reviewed above, it is evident that classical and quantum techniques complement each other. Under the Farm-Fetching framework, while classical models are effective in facilitating straightforward search queries, they cannot navigate intricate non-linear structural vagaries. In this regard, the novel QRAG model enhances semantic reasoning and retrieves generation through the inclusion of both frameworks into one structure.
